\section{From localized moments to a typical bridge}\label{sec:renewal}

Proposition~\ref{prop:08-localized-length} turns the confined first and
second moments into the following seed estimate: for every fixed
\(\xi>0\),
\[
 \Prob_h^{\rm br}\{|\gamma|\ge h^{4/3-\xi}\}\ge h^{-o(1)}.
\]
This probability may tend to zero. We amplify it to a high-probability
length bound by running independent irreducible-bridge trials, then
pay explicitly for conditioning on a specified terminal height.

The irreducible-bridge construction goes back to Kesten's renewal
identity~\cite{Kesten1963}; see Madras--Slade~\cite{MadrasSlade1993}
and the appendix of Lawler--Schramm--Werner~\cite{LSW2004}.
We derive the normalization and the conditional renewal estimates
needed for the present port convention.

A bridge is \emph{irreducible} if it has positive height and no
once-crossed intermediate level. Cutting at all such levels gives a unique
sequence of irreducible bridges. Conversely any sequence concatenates
freely: the pieces occupy successive disjoint open strips and their port
weights multiply. Let \(I_h\) be the total irreducible mass at height
\(h\ge1\), and set
\[
 I(z)=\sum_{h\ge1}I_hz^h,\qquad B(z)=\sum_{h\ge0}B_hz^h.
\]
Uniqueness gives \(B(z)=(1-I(z))^{-1}\) for \(0\le z<1\).
Since \(B_h\asymp(1+h)^{-1/4}\),
\[
 B(e^{-1/l})\asymp l^{3/4},\qquad
 \sum_{h\ge1}I_h=1,\qquad \sum_{h\ge l}I_h\le Cl^{-3/4}.
\]
For the tail, use
\((1-e^{-1})\sum_{h\ge l}I_h\le1-I(e^{-1/l})\); normalization follows
by letting \(z\uparrow1\).

We may therefore sample independent irreducible bridges, giving each
individual bridge probability equal to its critical weight. Their
concatenation is an infinite self-avoiding walk in the half-plane. Write
\(U_0=0<U_1<U_2<\cdots\) for its renewal heights, and let
\[
 T(l)=\min\{k:U_k>l\},\qquad G_l=\text{the prefix ending at }U_{T(l)}.
\]
The renewal heights tend to infinity because their increments are positive
integers. The event \(\mathcal H_h=\{h=U_k\text{ for some }k\}\) has probability
\(B_h\). Conditional on \(\mathcal H_h\), the prefix ending at height \(h\)
has exactly the critical height-\(h\) bridge law, normalized by \(B_h\).
This follows path by path from the unique irreducible decomposition.
The analogous strip-conditioning identity in the square-lattice vertex
convention appears in \cite[Proposition 2.10]{DyhrEtAl2011}; its final
bond lies above the conditioning height.

\begin{lemma}[Renewal overshoots and a successful trial]
\label{lem:08-renewal-trial}
For \(x\ge2l\ge2\),
\[
 \Prob\{U_{T(l)}>x\}\le C(l/x)^{3/4}.
\]
For every fixed \(\xi>0\), a trial beginning at a renewal and ending at
the first renewal more than \(h\) bands above it has probability at least
\(h^{-o(1)}\) of containing \(h^{4/3-\xi}\) vertices.
\end{lemma}
\begin{proof}
The last renewal before crossing \(l\) is some \(i\le l\), and its next
increment must exceed \(x-i\). Hence
\[
 \Prob\{U_{T(l)}>x\}
 =\sum_{i=0}^{l}B_i\sum_{j>x-i}I_j
 \le Cx^{-3/4}\sum_{i=0}^{l}B_i
 \le C(l/x)^{3/4}.
\]
For the second assertion let \(Z_h\) count renewals
\(m\in[h/2,h]\) whose prefixes have length at least
\(h^{4/3-\xi}\). Applying \eqref{eq:R-32d} with slack \(\xi/2\),
which absorbs the constant factor between \(m\) and \(h\), gives
\(\E Z_h\ge h^{3/4-o(1)}\). The unrestricted number \(N_h\) of
renewals up to \(h\) satisfies
\[
 \E N_h^2
 \le\sum_{m\le h}B_m+
       2\sum_{m<n\le h}B_mB_{n-m}
 \le Ch^{3/2}.
\]
The product in the second sum follows from independence after a renewal.
Thus \(\Prob\{Z_h>0\}\ge(\E Z_h)^2/\E N_h^2\ge h^{-o(1)}\).
Every prefix counted by \(Z_h\) lies in the trial.
\end{proof}

\begin{theorem}[Typical finite-bridge length]
\label{thm:08-bridge-length}
For every \(\varepsilon>0\) there is \(c_\varepsilon>0\) such that a
critical height-\(H\) bridge satisfies
\[
 \Prob_H^{\mathrm{br}}\{H^{4/3-\varepsilon}\le|\gamma|
                         \le H^{4/3+\varepsilon}\}
       \ge1-O_\varepsilon(H^{-c_\varepsilon}).
\]
Its diameter is at least \(d_0H\) and at most \(CH^{1+\varepsilon}\)
with polynomially small exceptional probability.
\end{theorem}
\begin{proof}
Fix small \(\theta,\xi>0\), and put
\[
 h=\lfloor H^{1-\theta}\rfloor,\qquad
 R=\lfloor H^{1-\theta/2}\rfloor,\qquad
 K=\lfloor H^{\theta/16}\rfloor.
\]
Run \(K\) successive trials of Lemma~\ref{lem:08-renewal-trial},
restarting at each completion. Independence of irreducible pieces after
these stopping renewal indices makes the trials independent. If
\(\Delta_j\) is their height advance, then
\[
 \Prob\left\{\sum_{j=1}^K\Delta_j>R\right\}
 \le\sum_{j=1}^K\Prob\{\Delta_j>R/K\}
 \le CK(hK/R)^{3/4}
 \le CH^{-17\theta/64}.
\]
For large \(H\) every trial succeeds with probability at least
\(H^{-\theta/32}\), by its \(h^{-o(1)}\) lower bound. Thus all trials
fail with probability at most \(\exp(-cH^{\theta/32})\). Outside these
two events, a successful trial is contained in \(G_R\). Choose
\(\theta,\xi\) so that
\((1-\theta)(4/3-\xi)>4/3-\varepsilon\). Then
\[
 \Prob\{|G_R|<H^{4/3-\varepsilon}\}
       \le CH^{-17\theta/64}+\exp(-cH^{\theta/32}).
\]

We transfer this estimate to conditioning on \(\mathcal H_H\). Put
\(\tau=T(R)\) and extend \(B_j\) by zero for \(j<0\). For any event
\(E\) determined by the prefix through \(\tau\), renewal independence
gives the exact identity
\[
 \Prob(E\mid \mathcal H_H)
    =\frac{\E[\mathbf1_E B_{H-U_\tau}]}{B_H}.
\]
Here \(R<H\), so a path which overshoots \(H\) at \(\tau\) could not
have hit \(H\) earlier. On \(U_\tau\le H/2\) the density ratio is
bounded by a constant. For all sufficiently large \(H\), we have \(R\le H/4\).
For the remaining conditional probability, let \(i\le R\) be the last
renewal before crossing \(R\), and let \(l\) be the height of the
crossing irreducible. If \(i+l>H/2\), then \(l\ge H/2-R\).
Summing over these possibilities gives the upper bound
\[
 \frac1{B_H}
 \sum_{l\ge H/2-R} I_l
       \sum_{i=0}^{\min(R,H-l)}B_iB_{H-l-i}
 \le C(R/H)^{1/2}.
\]
For completeness, writing
\(C_R(n)=\sum_{i=0}^{\min(R,n)}B_iB_{n-i}\), the full convolution
bound gives \(C_R(n)\le C(n+1)^{1/2}\le CR^{1/2}\) when
\(n\le2R\). When \(n>2R\), use
\(C_R(n)\le CR^{3/4}n^{-1/4}\le CR^{1/2}\).
Multiply this uniform bound by
\(\sum_{l\ge H/2-R}I_l\le CH^{-3/4}\) and
\(B_H^{-1}\le CH^{1/4}\). This proves the displayed estimate.
The prefix \(G_R\) is part of the conditioned height-\(H\) bridge
whenever its endpoint is at most \(H/2\), so the lower-length failure
is \(O(H^{-\theta/4})\). For example, for \(0<\varepsilon<1\),
\(\theta=\varepsilon/8\) and \(\xi=\varepsilon/4\) give the explicit
bound \(O_\varepsilon(H^{-\varepsilon/32})\).

The full strip first-length estimate in
Theorem~\ref{prop:02-length-reduction}, divided by \(B_H\), gives
\(\E_H^{\mathrm{br}}|\gamma|\le H^{4/3+o(1)}\). Markov's inequality proves the
upper bound with polynomially small failure. Finally, dividing the lateral
loss \(CHM^{-5/4}\) of Lemma~\ref{lem:08-localization} by \(B_H\) gives
\(\Prob_H^{\mathrm{br}}\{\text{lateral extent}>M\}\le C(H/M)^{5/4}\).
Take \(M=H^{1+\varepsilon}\). The deterministic height span is \(d_0H\),
which also proves the lower diameter bound.
\end{proof}
